\input{libraries/report-preamble.tex}
\title{A Bounded Binary64 Exponential}
\begin{document}
\maketitle

\section{Task and numerical domain}
\texttt{expTaylor6} approximates $e^x$ on $[-1,0]$ using the degree-six
Taylor polynomial
\[
  p_6(x)=1+x+\frac{x^2}{2}+\frac{x^3}{6}
         +\frac{x^4}{24}+\frac{x^5}{120}+\frac{x^6}{720}.
\]
The exponential power series appears in DLMF equation 4.2.19~\cite{series}.
On this interval its alternating tail has magnitude at most $1/5040$ in exact
arithmetic.  The target bound for the binary64 computation is $1/4000$.

\texttt{expTaylor6} accepts and returns binary64 bit patterns in
\texttt{UInt64} words.  The result is optional: a finite input in $[-1,0]$
returns a value, and other inputs return \texttt{none}.  Both signed zeros
return the exact bits of one.  Negative subnormals are admitted.  Infinities
and NaNs are rejected by the bit-range check.
LeanExe compiles the Lean function and its callers to WebAssembly (WASM).
The examples run that WASM in Wasmtime.

\section{Verification status}
Lean checks the source types.  Execution tests pass for selected inputs:
WASM and native Lean agree bit for bit, domain checks give the expected status,
and accepted results lie within $1/4000$ of JavaScript's \texttt{Math.exp}.

A separate Lean theorem, \texttt{Project.ExpSmall.Spec.expSmall\_real\_error},
proves termination and a finite positive result within $1/4000$ of the real
exponential for every finite binary64 input in $[-1,0]$~\cite{prior}.
Its subject is the WASM model \texttt{Project.ExpSmall.module}.
Formal verification of this library function and its compiled clients,
including their connection to that theorem, remains deferred.

\section{Implementation and coefficient choices}
\lstinputlisting{LeanExe/Lib/Transcendental/Exp/Basic.lean}

The listing comes from \texttt{Basic.lean}.  Each multiplication and addition
rounds separately.  The hexadecimal literals encode the rounded coefficients.
\texttt{expTaylor6Polynomial} exposes the unguarded evaluator.
\texttt{expTaylor6} applies the domain guard before calling it.
Evaluation uses six multiplications and six additions.

The coefficient words, evaluation order, and domain follow
\texttt{Project.ExpSmall.Model}.  This component expresses the arithmetic
through \texttt{LeanExe.Float64} primitives.

\section{A decay client}
The decay client approximates the remaining fraction $e^{-t}$ for dimensionless
time $t\in[0,1]$.  It checks $t$, negates its binary64 sign bit, and calls the
exponential component.
The commands run from the repository root after following the
\href{https://github.com/jsmorph/leanexe/blob/lib1/DEVELOPING.md}{setup guide}.

\begin{lstlisting}
tools/seminum exp-taylor6 -0.5
0.6065321180555556
tools/seminum decay 0.5
0.6065321180555556
tools/seminum exp-taylor6-bits bfe0000000000000
3fe368b60b60b60c
\end{lstlisting}

The host converts decimal inputs to binary64 and formats returned values in
decimal.  The hexadecimal command preserves input and output bits.
Out-of-domain inputs cause exit status 2 and a diagnostic on standard error.

\section{Tests}
From the repository root, \texttt{node test/seminum.js} runs the comparisons
described above at 33 equally spaced points in $[-1,0]$, both signed zeros,
subnormals, adjacent boundary words, infinities, and NaNs.  Accepted outputs
are finite and positive on these cases.  The largest sampled absolute
difference from \texttt{Math.exp} is about $1.7611438411335723\times10^{-4}$
at $x=-1$, within the target $1/4000$.  The decay client has corresponding
comparisons on $[0,1]$.  CLI tests check decimal and hexadecimal output and
rejected input.  The cases are defined in \texttt{test/seminum.js} and the
native reference driver \texttt{test/SeminumNative.lean}.

\begin{thebibliography}{9}
\bibitem{series} NIST Digital Library of Mathematical Functions, equation 4.2.19.
\url{https://dlmf.nist.gov/4.2.E19}.
The power series supplies the coefficients before rounding and truncation.
\bibitem{prior}
\href{https://github.com/jsmorph/leanexe/blob/lib1/proofs/talos/lean/Project/ExpSmall/Spec.lean}{LeanExe, \emph{Small-exponential execution theorem}.}
The numerical proof bounds polynomial approximation, coefficient rounding,
and arithmetic rounding for \texttt{Project.ExpSmall.module}.
\end{thebibliography}
\end{document}
